\subsection{Parameter-Free Realization of Unit Circle Area Phase Compression: Cayley Compactification and Sector Area Coordinates}\label{subsec:unit-circle-phase-lift}
For any angular variable \(\vartheta\in(-\pi,\pi)\), the corresponding sector area of the unit circle is \(\vartheta/2\) and the full circle area is \(\pi\), so the sector area fraction is \(\vartheta/(2\pi)\in(-\tfrac12,\tfrac12)\).
Accordingly, for \(x\in\RR\) we define the \textbf{area phase coordinate} (parameter-free compactification)
\begin{equation}\label{eq:unit-circle-phase-upsilon}
\boxed{\ \upsilon(x):=\frac{1}{\pi}\arctan(x)\in\Bigl(-\tfrac12,\tfrac12\Bigr)\ }.
\end{equation}
and set
\begin{equation}\label{eq:unit-circle-phase-vartheta}
\boxed{\ \vartheta(x):=2\pi\,\upsilon(x)=2\arctan(x)\in(-\pi,\pi)\ }.
\end{equation}
Then \(\upsilon(x)=\vartheta(x)/(2\pi)\) may be interpreted as the ``signed sector area fraction from \(0\) to \(x\).''

\paragraph{Unit Circle Embedding (Cayley Form)}
Define
\begin{equation}\label{eq:unit-circle-cayley-embed}
\boxed{\ \iota(x):=e^{\mathrm{i}\vartheta(x)}=e^{2\mathrm{i}\arctan(x)}=\frac{1+\mathrm{i}x}{1-\mathrm{i}x}\in\TT\ }.
\end{equation}
This yields an analytic homeomorphism from \(\RR\) to \(\TT\setminus\{-1\}\) satisfying the basepoint normalization \(\iota(0)=1\).
As \(x\to\pm\infty\) one has \(\vartheta(x)\to\pm\pi\) and \(\iota(x)\to-1\), thereby compressing ``infinity'' to a single point on the unit circle boundary.
If we set \(\widehat{\RR}:=\RR\cup\{\infty\}\) as a one-dimensional projective compactification and adopt the convention \(\iota(\infty):=-1\), then \(\iota:\widehat{\RR}\to\TT\) becomes a homeomorphism; in this sense, the ``unbounded variable'' is parameter-free compressed to a single boundary endpoint on the unit circle.
The above M\"obius homeomorphism and boundary extension are standard facts in complex analysis (cf.\ \cite{Ahlfors1979ComplexAnalysis}).
The explicit inverse projection is
\begin{equation}\label{eq:unit-circle-phase-inverse}
\boxed{\ \upsilon^{-1}(u)=\tan(\pi u),\qquad x=\tan\!\Bigl(\frac{\vartheta}{2}\Bigr)=\tan(\pi u),\qquad
x=\mathrm{i}\,\frac{1-w}{1+w}=-\mathrm{i}\,\frac{w-1}{w+1}\ (w\in\TT,\ w\neq-1)\ }.
\end{equation}

\paragraph{Push-Forward Uniformization of the Cauchy Weight (Pullback of the Unit Circle Haar Measure)}
From \(\vartheta(x)=2\arctan(x)\) one directly computes
\[
\frac{\dd\vartheta}{\dd x}=\frac{2}{1+x^{2}},
\qquad\Longrightarrow\qquad
\boxed{\ \frac{\dd\vartheta}{2\pi}=\dd\upsilon=\frac{1}{\pi}\frac{\dd x}{1+x^{2}}\ }.
\]
Therefore, integration on the real axis weighted by the Cauchy kernel \(\frac{1}{\pi(1+x^{2})}\) is strictly equivalent to uniform circular averaging on the unit circle; ``near-zero weighting and far-end suppression'' is not an externally imposed rule but rather a structural consequence uniquely induced by the Jacobian of the compactification gate.

\begin{corollary}[Cayley--Haar compatibility forces the Cauchy weight (unique normalization)]\label{cor:unit-circle-cayley-haar-forces-cauchy}
Let the standard Cayley transform (real axis \(\leftrightarrow\) unit circle) be
\[
\mathcal{C}(x):=\frac{x-\mathrm{i}}{x+\mathrm{i}},
\qquad
\mathcal{C}^{-1}(y)=\mathrm{i}\,\frac{1+y}{1-y},
\]
and let \(\iota\) be the phase embedding of \eqref{eq:unit-circle-cayley-embed}.
Then \(\mathcal{C}(x)=-\iota(x)\), so the Cayley transport and the phase gate of this appendix are the same coordinate system modulo M\"obius equivalence (differing only by a fixed rotation on the unit circle).
In particular, the pullback of the normalized Haar measure \(\dd\vartheta/(2\pi)\) on the unit circle under this coordinate must be
\[
\dd\mu(x)=\frac{1}{\pi(1+x^{2})}\,\dd x,
\]
and under the discipline ``preserve Haar normalization and compress \(\infty\) to a single point,'' this probability weight is uniquely determined by the Jacobian.
\end{corollary}
\begin{proof}
From \eqref{eq:unit-circle-cayley-embed},
\[
\iota(x)=\frac{1+\mathrm{i}x}{1-\mathrm{i}x}
\quad\Longrightarrow\quad
-\iota(x)=\frac{-1-\mathrm{i}x}{1-\mathrm{i}x}
=\frac{x-\mathrm{i}}{x+\mathrm{i}}=\mathcal{C}(x),
\]
so the two coordinates differ only by a rotation factor \(-1\).
The Haar pullback formula is simply the Jacobian computation \(\dd\vartheta/(2\pi)=\dd x/(\pi(1+x^2))\), yielding the displayed weight.
Uniqueness follows from the fact that given the Haar normalization and \(\mathcal{C}\) (or \(\iota\)) as a M\"obius homeomorphism on the projective compactification, the pullback measure is uniquely determined by the change-of-variables formula. \qed
\end{proof}

\paragraph{Phase Addition on the Real Line: M\"obius Group Law on the Projective Compactification (Overflow Immune)}
Write the projective compactification as \(\overline{\RR}:=\RR\cup\{\infty\}\) and extend \(\iota\) to the limit \(x\to\infty\) by \(\iota(\infty):=-1\).
\begin{proposition}[M\"obius group law for the half-angle tangent parameter]\label{prop:unit-circle-phase-mobius-law}
On \(\overline{\RR}\) define the binary operation
\begin{equation}\label{eq:unit-circle-phase-mobius-law}
\boxed{\;
x\oplus y:=
\begin{cases}
\dfrac{x+y}{1-xy}, & x,y\in\RR,\ 1-xy\neq 0,\\[6pt]
\infty, & x,y\in\RR,\ 1-xy=0,
\end{cases}
}
\end{equation}
and adopt the continuous extension convention
\(
x\oplus \infty=\infty\oplus x:=-\frac{1}{x}
\)
(with \(1/0=\infty,\ 1/\infty=0\)).
Then \((\overline{\RR},\oplus)\) is a commutative group with identity element \(0\) and inversion \(x\mapsto -x\), and \(\iota:(\overline{\RR},\oplus)\to(\TT,\cdot)\) is a group homomorphism:
\begin{equation}\label{eq:unit-circle-phase-mobius-hom}
\boxed{\ \iota(x\oplus y)=\iota(x)\cdot\iota(y)\ }.
\end{equation}
In particular, \(\infty\) is a legitimate boundary element in this group and satisfies
\begin{equation}\label{eq:unit-circle-phase-inversion-as-translation}
\boxed{\ x\oplus\infty=-\frac1x\ },
\end{equation}
so that ``inversion'' can be internalized as a single group operation with the boundary element.
\end{proposition}
\begin{proof}
For \(x\in\RR\) write \(\vartheta(x)=2\arctan(x)\).
By the trigonometric identity
\(
\tan\!\bigl(\frac{\alpha+\beta}{2}\bigr)=\frac{\tan(\alpha/2)+\tan(\beta/2)}{1-\tan(\alpha/2)\tan(\beta/2)}
\)
with \(\alpha=\vartheta(x),\beta=\vartheta(y)\), one obtains for \(x,y\in\RR\) with \(1-xy\neq 0\)
\(
\tan\!\bigl((\vartheta(x)+\vartheta(y))/2\bigr)=\frac{x+y}{1-xy}
\), which is \eqref{eq:unit-circle-phase-mobius-law}.
When \(1-xy=0\), the quantity \((\vartheta(x)+\vartheta(y))/2\) hits the principal value \(\pm\pi/2\), yielding \(x\oplus y=\infty\).
The case involving \(\infty\) follows by taking the limit \(y\to\infty\), giving \(x\oplus\infty=-1/x\) and \(\infty\oplus\infty=0\).
Finally, from
\(
\iota(x)=e^{\mathrm{i}\vartheta(x)}
\)
one deduces
\(
\iota(x)\iota(y)=e^{\mathrm{i}(\vartheta(x)+\vartheta(y))}=e^{\mathrm{i}\vartheta(x\oplus y)}=\iota(x\oplus y)
\), yielding \eqref{eq:unit-circle-phase-mobius-hom}; commutativity and associativity are inherited from those of multiplication on the unit circle. \qed
\end{proof}

\begin{remark}[The half-angle phase term in completion coincides with the phase arithmetic coordinate]\label{rem:unit-circle-phase-completion-half-angle}
If the completion/unit-circle twist channel admits a strict phase--amplitude separation (Proposition~\ref{prop:completion-endpoint-jet-diffusion}), i.e., for \(u=e^{it}\) one has
\(
P(it)=\frac{it}{2}+g(S(t)),\ S(t)=2\cos(t/2)
\),
then rewriting in the half-angle tangent parameter \(x=\tan(t/2)\) (equation~\eqref{eq:unit-circle-phase-inverse}) yields
\[
\boxed{\;
P\!\bigl(it(x)\bigr)
=\mathrm{i}\arctan(x)\ +\ g\!\left(\frac{2}{\sqrt{1+x^2}}\right)\;},
\qquad t(x)=2\arctan(x).
\]
Thus the ``half-angle phase term'' \(\frac{it}{2}\) forced by the completion is strictly identical to the phase arithmetic variable (\(\frac{it}{2}=\mathrm{i}\arctan x\)); the amplitude enters only through the invariant \((1+x^2)^{-1/2}\).
\end{remark}
